%==============================================================================
% Capitulo 10: Recomendaciones y lineas de trabajo futuro
%==============================================================================
\section{Recomendaciones y lineas de trabajo futuro}

A partir de los resultados obtenidos, del analisis de errores presentado en la Seccion~7.4 y de las limitaciones reconocidas en la Seccion~8.3, se formulan las siguientes recomendaciones para la evolucion del sistema y para futuras replicaciones del estudio.

\subsection{Incorporacion de un clasificador supervisado con corpus propio}

Se sugiere incorporar progresivamente un clasificador supervisado entrenado con el corpus etiquetado generado por el propio sistema durante su operacion. La acumulacion sostenida de decisiones validadas por los sectores responsables constituye un activo informacional valioso que, una vez alcanzado un volumen del orden de 5.000 casos, permitiria entrenar modelos de clasificacion especificos del dominio organizacional. Las ventajas potenciales de esta migracion incluyen: reduccion del costo marginal por incidente clasificado al eliminar la dependencia del servicio externo de inferencia, disminucion de la latencia al ejecutar la inferencia localmente, y ganancia de soberania sobre los datos al evitar la transferencia internacional de fragmentos textuales. \textcite{Touvron2023_llama2} y trabajos subsiguientes sobre modelos de codigo abierto sugieren que los LLM desplegados localmente alcanzan rendimientos competitivos en tareas de clasificacion cuando se invocan mediante prompts adecuadamente disenados, lo que abre una via concreta hacia un sistema completamente autoalojado.

\subsection{Ampliacion de canales de entrada}

Se propone ampliar el conjunto de canales de entrada para incluir mensajeria corporativa instantanea, particularmente Slack y Microsoft Teams en aquellas organizaciones donde estas plataformas constituyen el canal de comunicacion dominante. La inclusion de aplicaciones moviles propias para usuarios de campo constituye una extension natural que atiende la demanda creciente de canales asincronicos integrados al ecosistema laboral cotidiano. La arquitectura modular propuesta en el Capitulo~5 facilita esta ampliacion: cada nuevo canal requiere unicamente un disparador adicional en N8N, sin modificaciones en las capas de procesamiento ni de persistencia.

\subsection{Panel de monitoreo en tiempo real}

Se sugiere implementar un panel de monitoreo sustentado sobre tecnologias de observabilidad maduras, tales como Prometheus para la captura de metricas y Grafana para la visualizacion. Este panel permitiria supervisar la latencia del sistema, la tasa de aciertos del clasificador, la distribucion de carga entre sectores y el costo agregado de las invocaciones al modelo externo, habilitando ajustes operativos basados en evidencia continua y alertas tempranas frente a degradaciones del servicio. La existencia de endpoints de salud (\texttt{GET /api/v1/health}) en cada componente del sistema proporciona los puntos de datos necesarios para instrumentar esta capa de observabilidad.

\subsection{Mecanismos de aprendizaje activo}

Se propone incorporar mecanismos de aprendizaje activo (\emph{active learning}), en virtud de los cuales aquellos casos clasificados con baja confianza se prioricen automaticamente para revision humana, y los resultados de esa revision retroalimenten al sistema mediante un mecanismo de ajuste incremental. Este enfoque maximiza el rendimiento marginal de cada intervencion humana ---que de todos modos ocurre en el 9,5~\% de los casos--- y acelera la mejora continua del clasificador con un costo operativo controlado.

\subsection{Estudios de replicacion}

Se sugiere realizar estudios de replicacion en organizaciones de distinto tamano, sector y region geografica, con el objeto de fortalecer la validez externa de los hallazgos y construir un cuerpo de evidencia empirica sobre el desempeno de soluciones similares en contextos heterogeneos. Particularmente relevantes serian las replicas en organizaciones del sector publico (con mayor proporcion de incidentes vinculados a sistemas administrativos), en organizaciones del sector industrial (con mayor proporcion de incidentes vinculados a infraestructura fisica) y en organizaciones de mayor tamano donde el volumen diario de incidentes supere las 100 unidades.

\subsection{Evaluacion de modelos de lenguaje de codigo abierto}

Se propone evaluar la integracion de modelos de lenguaje de codigo abierto desplegados localmente ---tales como Llama~3, Mistral o Qwen--- como alternativa al modelo externo Gemini~2.5 Flash. Esta evaluacion permitiria cuantificar el costo en exactitud frente a la ganancia en soberania de datos, la eliminacion de la dependencia de un proveedor unico y la reduccion de costos operativos a largo plazo. \textcite{Touvron2023_llama2} sugieren que los modelos abiertos contemporaneos alcanzan rendimientos competitivos en tareas de clasificacion, y la realizacion de un benchmark comparativo sobre el mismo corpus de 200 casos constituiria una contribucion adicional significativa al estado del arte.

\newpage
